% Section on Pocket Dosimeters
\section*{Dosímetros de Bolsillo}

Los dosímetros de bolsillo son una pequeña cámara de ionización, tipo condensador, de forma cilíndrica llena de aire y provista de un electrodo central.

Estos dispositivos deben cargarse antes de exponerlos a la radiación; esto se logra por medio de un cargador que les aplica un voltaje. Al exponer el dosímetro a la radiación, el aire de la cámara se ioniza y los iones se mueven hacia los electrodos, lo que ocasiona una pérdida de carga, la cual es proporcional a la cantidad de radiación recibida.

Los dosímetros de bolsillo pueden ser de lectura directa o de lectura indirecta. Los primeros están provistos de una fibra de cuarzo y de un pequeño microscopio, el cual permite observar el desplazamiento de fibra a lo largo de una escala graduada. Este desplazamiento se debe a que al perder carga el electrodo atrae a la fibra de cuarzo; por lo tanto el movimiento de la imagen de esta es una medida de la cantidad de radiación. En los de lectura indirecta, esta se efectúa mediante un dispositivo lector-cargador, el cual contiene un electrometro que permite estimar la exposición a la dosis midiendo la carga perdida por el citado dosímetro. Estos dosímetros se utilizan para determinar dosis diarias.

Los dosímetros de bolsillo están sujetos a descargas cuando se tiran o golpean, por lo que debe tenerse cuidado en su manejo. Cuando los dosímetro son expuestos a la humedad, este almacenamiento puede perderse con el consiguiente aumento de la corriente de fuga. La rapidez de fuga normal en un buen dosímetro debe ser del orden de 3\% en un período de 48 hrs.

% Section on Film Dosimeters
\section*{Dosímetros de Película}

Los dosímetros de película, consisten de un paquete de 2 o 3 placas fotográficas con plata (para rayos X, gamma y neutrones) protegidas de la luz y colocadas en un chasis. Estas películas constan de una base de acetato de celulosa o de algún plástico, revestida por ambos lados de una capa gelatinosa fotosensible (emulsión). De un lado se usa una emulsión muy sensible, mientras que del otro la emulsión es prácticamente insensible.

Este tipo de dosímetros se basa en el hecho de que la acción de la radiación ionizante o de sus emisiones secundarias producen un conglomerado microscópico de átomos de plata que constituyen la imagen latente, la cual al ser revelada produce un oscurecimiento visible de la emulsión.
